\documentclass{rescience}
\usepackage{booktabs}

\begin{document}
\begin{tikzpicture}[remember picture, overlay]
  \node[draw=darkred, fill=darkred, text=white, anchor=north west,
        minimum width=3.75cm, minimum height=10mm, yshift=-34mm]
        at (current page.north west) {\small \SpaceRoboto RESCIENCE C};
\end{tikzpicture}
{\let\newpage\relax\maketitle}
\marginnote{\footnotesize\sffamily
  \textbf{Received}\\\dateRECEIVED\\~\\
  \textbf{Code}\\\href{\codeURL}{\detokenize\expandafter{\codeURL}}\\~\\
  \textbf{Contact}\\\href{mailto:\contactEMAIL}{\contactEMAIL}}



\section{Introduction}
Computational results are most useful when another researcher can inspect the
path from input to claim. This starter presents a small reproducibility study
using a deterministic simulation as its running example. The design follows
the ReScience convention of keeping article metadata separate from article
content and makes the FAIR principles a concrete checklist~\parencite{wilkinson2016}.

\section{Methods}
We simulated a discrete-time population with carrying capacity $K=100$ and
growth rate $r=0.12$. At each step, the population was updated with
\begin{equation}
  N_{t+1}=N_t+rN_t\left(1-\frac{N_t}{K}\right).
  \label{eq:logistic}
\end{equation}
The exact parameters, code location, and data location should be replaced by
the author. Keeping them beside the manuscript makes a replication attempt
auditable rather than dependent on undocumented defaults.

\section{Results}
Starting from $N_0=12$, Eq.~\eqref{eq:logistic} approaches the carrying
capacity without overshoot. A replication should compare the full time series,
not only its final value, because two implementations can agree at the end
while differing substantially in the transient phase.

\begin{table}[ht]
\centering
\caption{Suggested replication checkpoints.}
\begin{tabular}{lll}
\toprule
Checkpoint & Record & Acceptance rule \\
\midrule
Input & parameters and initial state & exact match \\
Process & software and random seed & documented \\
Output & final value and full series & tolerance stated \\
\bottomrule
\end{tabular}
\end{table}

\section{Discussion}
The example is deterministic, so a mismatch points to an implementation or
parameter difference rather than sampling noise. For stochastic experiments,
the same section should report the seed policy, number of repetitions, and
uncertainty summary.

\section{Conclusion}
Use this file with the accompanying `metadata.tex` as a compact starting point
for a ReScience C submission. Replace the example claims with a replication of
an actual computational result and keep the metadata synchronized with the
PDF.

\hypersetup{linkcolor=black,urlcolor=darkgray}
\renewcommand\emph[1]{{\bfseries #1}}
\setlength\bibitemsep{0pt}
\printbibliography

\end{document}
